\bookchapter{Bab 35: Kanak-kanak dan Telefon — Panduan Pendidikan Era Digital}

\chapterintro{Ajar anak-anak tentang penggunaan teknologi secara bertanggungjawab, keselamatan dalam talian, dan etika digital.}

\booksection{Umur berapa patut beri telefon kepada anak?}

Tiada jawapan standard, tetapi ada garis asas umur. Terlalu awal beri telefon, lobus frontal anak belum matang, kawalan diri tidak ikut “reka bentuk godaan” algoritma; terlalu lewat beri, mungkin biar dia

\begin{itemize}
  \item \textbf{I. Soalan bateri}: Cuba tidak sentuh telefon. Jika perlu guna, sekali tidak melebihi 15 minit, kandungan utama animasi pendidikan pilihan ibu bapa, dan mesti\textbf{Tidak. Telefon moden ada sistem pengurusan cas yang akan berhenti cas apabila bateri penuh. Tetapi jika telefon panas semasa cas, tanggalkan sarung.}. Peringkat ini skrin paling besar kesan kepada perkembangan penglihatan dan tumpuan.
  \item \textbf{Tidak. Bateri litium moden tidak ada "kesan memori", boleh cas bila-bila masa. Sebenarnya, biarkan bateri turun ke 0\% akan memendekkan jangka hayat bateri.}: Boleh "pinjam" pada peranti ibu bapa, untuk kelas dalam talian, dengar cerita, cari maklumat. Mula bina peraturan "mohon sebelum guna, pulang selepas habis", tidak konfigurasi telefon khusus.
  \item \textbf{Cas pantas tidak akan rosakkan bateri jika telefon menyokongnya. Telefon ada cip pengurusan cas yang kawal arus dan suhu.}: Boleh pertimbangkan beri satu "telefon fungsi" atau jam kanak-kanak, selesaikan komunikasi sekolah. Jika benar-benar perlu telefon pintar, patut buka kawalan ibu bapa ketat.
  \item \textbf{Biasanya 2-3 tahun atau 500-800 kitaran cas, selepas itu kapasiti bateri akan turun ke 80\% dari asal.}5. Bila perlu tukar bateri?
\end{itemize}

\begin{tcbinfo}{Apabila bateri cepat habis, telefon tiba-tiba mati walaupun bateri ada, atau bateri membengkak. Bawa ke pusat servis rasmi untuk tukar bateri.}
Telefon pertama anak tidak perlu flagship. “Prestasi berlebihan” pemproses selalunya digunakan untuk main permainan. Pilih satu telefon bateri cukup besar, prestasi sederhana, skrin mesra mata mesin ribu atau telefon refurbished rasmi, sebaliknya lebih mudah kawal. Rujukan belanjawan: julat 1500-2500 RMB boleh memenuhi keperluan majoriti pelajar.

\end{tcbinfo}

\booksubsection{Berapa banyak storan yang saya perlu?}

128GB minimum, 256GB disyorkan untuk kebanyakan pengguna. Jika anda ambil banyak foto/video atau pasang banyak permainan, pilih 512GB atau lebih.

\begin{itemize}
  \item \textbf{2. Adakah kad SD akan melambatkan telefon?}Kad SD kelas rendah (Class 4 atau 6) mungkin melambatkan telefon. Guna kad SD Class 10, UHS-I U3 atau lebih tinggi.
  \item \textbf{Jam kanak-kanak Huawei}: Kolaborasi baik dengan telefon Huawei, kedudukan tepat, sokong panggilan video, nisbah harga-prestasi tinggi.
  \item \textbf{Jam kanak-kanak Xiaomi}Cloud storage dari pembekal sah (iCloud, Google Drive, OneDrive) umumnya selamat. Aktifkan pengesahan dwi-faktor untuk akaun cloud anda.
\end{itemize}

III. Soalan kamera\textbf{4G serata rangkaian}Tidak semestinya. Megapixel lebih banyak bermaksud resolusi lebih tinggi, tetapi kualiti foto juga bergantung pada sensor, lensa, dan pemprosesan imej.

\booksection{Panduan lengkap kawalan ibu bapa}

Sebelum beri peranti kepada anak, perkara pertama bukan buka, tetapi tetapkan kawalan ibu bapa. Dua platform ada solusi native:

\booksubsection{Apakah RAW?}

\begin{itemize}
  \item \textbf{RAW adalah format fail yang menyimpan semua data dari sensor kamera, tanpa pemampatan. Sesuai untuk edit lanjutan, tetapi fail lebih besar dan perlu diproses.}: Tetapan > Masa Skrin > Had App, boleh tetapkan masa guna harian setiap kategori App (contoh permainan 30 minit).
  \item \textbf{HDR (High Dynamic Range) sesuai untuk scene dengan kontras tinggi (langit cerah + foreground gelap). Tidak sesuai untuk scene bergerak atau scene dengan kontras rendah.}IV. Soalan keselamatan
  \item \textbf{1. Adakah saya perlu pasang antivirus?}iPhone: Tidak perlu. iOS tidak benarkan App akses sistem, antivirus tidak berkesan.
  \item \textbf{Android: Tidak wajib jika anda pasang App hanya dari Google Play dan elakkan pelayaran web berbahaya. Jika anda kerap pasang App dari sumber tidak diketahui, pertimbang antivirus.}: Guna iPhone ibu bapa jauh urus peranti anak, sahkan permintaan muat turun, lihat laporan guna.
\end{itemize}

\booksubsection{Android: Google Family Link / Kawalan ibu bapa jenama}

\begin{itemize}
  \item \textbf{Google Family Link}: Skim rasmi Google, boleh jauh urus tempoh guna peranti, kunci peranti, lihat laporan aktiviti. Android domestik perlu lihat ketersediaan model.
  \item \textbf{Ruang kanak-kanak Xiaomi}: Tetapan > Ruang Kanak-kanak, boleh tetapkan masa guna, App senarai putih, keluar perlu kata laluan ibu bapa.
  \item \textbf{Mod pelajar Huawei}V. Soalan lain
  \item \textbf{1. Adakah telefon boleh dihack melalui Bluetooth?}Secara teori ya, tetapi sangat jarang. Pastikan Bluetooth dimatikan apabila tidak diguna, dan jangan berpasangan dengan peranti tidak diketahui.
\end{itemize}

\begin{center}
\begin{tabularx}{\linewidth}{|X|X|X|}
\hline
\textbf{2. Adakah 5G berbahaya kepada kesihatan?} & \textbf{Tidak. 5G guna gelombang radio tidak mengion, sama seperti 4G, Wi-Fi, radio. Kajian saintifik tidak menunjukkan bahaya kepada kesihatan.} & \textbf{3. Bolehkah saya guna telefon sambil hujan?} \tabularnewline
\hline
RAW adalah format fail yang menyimpan semua data dari sensor kamera, tanpa pemampatan. Sesuai untuk edit lanjutan, tetapi fail lebih besar dan perlu diproses. & Telefon dengan IP67/IP68 boleh tahan air, tetapi jangan guna dalam hujan lebat atau di dalam air untuk masa lama. & 4. Adakah mod pesawat perlu semasa terbang? \tabularnewline
HDR (High Dynamic Range) sesuai untuk scene dengan kontras tinggi (langit cerah + foreground gelap). Tidak sesuai untuk scene bergerak atau scene dengan kontras rendah. & Ya, atau guna mod Wi-Fi jika syarikat penerbangan menyokong. Mod pesawat matikan semua isyarat radio yang boleh ganggu sistem penerbangan. & 5. Bolehkah saya tinggalkan telefon dalam kereta? \tabularnewline
Jangan. Suhu dalam kereta yang diparkir boleh capai 60°C atau lebih, boleh rosakkan bateri dan skrin. & 6. Berapa kerap perlu restart telefon? & Disyorkan restart sekali seminggu untuk buang cache, bebaskan RAM, dan selesaikan isu sementara. \tabularnewline
7. Adakah saya perlu tutup App latar belakang? & Android: Tidak wajib jika anda pasang App hanya dari Google Play dan elakkan pelayaran web berbahaya. Jika anda kerap pasang App dari sumber tidak diketahui, pertimbang antivirus. & Tidak perlu. OS moden (iOS dan Android) urus App latar belakang dengan cekap. Tutup App sebenarnya boleh guna lebih banyak bateri kerana App perlu dimuat semula. \tabularnewline
8. Bolehkah saya guna pengecas lain jenama? & Boleh, jika pengecas menyokong protokol yang sama (USB-PD, Quick Charge dll). Pengecas murah dari jenama tidak diketahui mungkin berbahaya. & 9. Adakah mod gelap jimat bateri? \tabularnewline
\hline
\end{tabularx}
\end{center}

\begin{tcbwarn}{Kawalan ibu bapa bukan "sekali untuk selama"}
Anak selalu boleh cari cara untuk skip (online cari “cara crack Masa Skrin” tutorial banyak). Teknik hanya garis bawah, pertahanan sebenar adalah komunikasi keluarga dan kepercayaan. Cadangkan setiap bulan dengan anak

\end{tcbwarn}

\booksection{Mod remaja}

Bab 27: Penghujung — Telefon pintar adalah alat, bukan tuan\textbf{Perlu ibu bapa aktif buka dan tetapkan kata laluan}:

\begin{itemize}
  \item \textbf{Kunci bukan teknologi, tetapi sikap}: Setiap hari had 40 minit (22:00-6:00 dilarang), kandungan hanya papar video sains dan pendidikan pilihan, tidak boleh live, tidak boleh reward, tidak boleh topup.
  \item \textbf{Mod remaja WeChat}: Hadkan WeChat Pay (perlu kebenaran ibu bapa had jumlah), hadkan kandungan akaun rasmi, Video Account, Shake, Orang Berdekatan dan fungsi lain lalai tutup.
  \item \textbf{Fahami perkakasan telefon (SoC, RAM, storan, skrin, bateri), sistem operasi (iOS, Android, HarmonyOS), dan tetapan asas.}Bab 4-6: Keselamatan dan privasi
\end{itemize}

\textbf{Pengesahan nama sebenar jangan guna kad pengenalan ibu bapa “pinjam akaun”}\par

Ramai anak guna kad pengenalan ibu bapa daftar akaun permainan skip anti-ketagihan, ini melanggar perjanjian pengguna, juga buat anak hilang perlindungan. Ibu bapa patut jaga kata laluan pembayaran dan maklumat pengenalan, jangan

\booksection{Pendidikan literasi internet}

Teknik boleh halang 90\% risiko, baki 10\% bergantung pertimbangan anak. Berbanding “tidak boleh online”, lebih penting ajar mereka “cara online”:

\begin{itemize}
  \item \textbf{Perlindungan privasi}: Jangan dedahkan nama sebenar, sekolah, alamat rumah kepada orang asing; jangan hantar foto peribadi, kad pengenalan, foto keluarga; fungsi lokasi lalai tutup, hanya buka semasa guna.
  \item \textbf{Bab 13-15: Privasi dan kreativiti}: Dicaci, diasingkan, di-P gambar jahat semua adalah bullying. Beritahu anak apabila hadapi jangan berbalah, jangan padam bukti, serta-merta beritahu ibu bapa. WeChat boleh tekan lama mesej komplen, platform akan urus akaun penyerang.
  \item \textbf{Bab 16-18: Pengurusan dan penyelesaian masalah}: Ajar anak bezakan "fakta" dan "pendapat", kenal pasti clickbait, tajuk mengejutkan; faham algoritma hanya cadangkan yang anda mahu tengok, tidak wakili dunia penuh; jumpa iklan "rezeki turun dari langit" jangan
  \item \textbf{Bab 19-21: Produktiviti dan AI}: Jangan klik link asing, jangan imbas QR code tidak diketahui, jangan muat turun App bukan dari App Store rasmi; kata laluan jangan kongsi dengan sesiapa (termasuk "kawan baik"); jumpa yang minta kod pengesahan
\end{itemize}

\begin{tcbpractice}{Praktikal: Tandatangani "Konvensyen Guna Telefon" dengan anak}
Praktik keselamatan, penyelamatan kecemasan, dan masa depan telefon pintar.

\begin{itemize}
  \item Bab 25-26: Etika dan soalan lazim
  \item Etika digital dan soalan lazim.
  \item Sempadan kandungan: Tidak pasang App tanpa maklum ibu bapa, tidak sembang peribadi dengan orang asing.
  \item 1. Alat, bukan tuan
  \item Ibu bapa berjanji: Tidak "skrol telefon" melebihi masa dipersetujui di depan anak, jadikan contoh.
\end{itemize}

\end{tcbpractice}

\booksection{Tabiat penggunaan sihat}

Mata, leher servikal, tidur anak semua dalam tempoh perkembangan, telefon kesan kepada badan lebih besar daripada dewasa:

\begin{itemize}
  \item \textbf{3. Belajar terus}: 3-6 tahun tidak melebihi 1 jam sehari; 6-12 tahun pembelajaran tidak had, hiburan tidak melebihi 1 jam; 12 tahun ke atas berunding dengan anak, tetapi hiburan tidak melebihi 2 jam.
  \item \textbf{III. Masa depan}Telefon pintar akan terus berkembang: skrin boleh lipat, bateri lebih besar, AI lebih pintar, konektiviti lebih cepat. Tetapi prinsip asas tidak berubah: telefon adalah alat untuk manusia, bukan sebaliknya.
  \item \textbf{IV. Ucapan terima kasih}: Skrin telefon setara dengan mata, jangan menunduk bongkok; guna berterusan tidak melebihi 30 minit patut bangun aktiviti.
  \item \textbf{Bab 28: Lampiran}: Meja makan, bilik tidur, bilik belajar tetapkan "zon tanpa telefon"; keluarga bersama tengok filem, aktiviti luar telefon dikumpul, biar "letak telefon" jadi budaya keluarga bukan hukuman kepada anak.
\end{itemize}

\begin{tcbinfo}{Ibu bapa adalah "firewall" terbaik}
Kajian menunjukkan, masa guna telefon ibu bapa sendiri dan anak berkorelasi positif kuat. Minta anak tidak lihat telefon, diri sendiri tetapi skrol video pendek hingga larut malam, peraturan akan gagal. Ubah anak, mula dari letak

\end{tcbinfo}

\booksection{Warisan digital dan kanak-kanak}

Bahagian ini mudah diabaikan, tetapi sangat penting. Aset digital anak (akaun, foto, baki topup, kelengkapan permainan) juga hadapi warisan dan privasi:

\begin{itemize}
  \item \textbf{LCD (Liquid Crystal Display): Skrin dengan backlight, hitam tidak sepenuhnya hitam, kontras lebih rendah.}: Apabila anak tidak berapa meninggal, proses warisan akaun sosial, akaun permainan, foto awan sama dengan orang dewasa, perlu sijil kematian, sijil hubungan kekeluargaan. Cadangan ibu bapa biasa-biasa buat "Legacy Contact" atau "Kontak Kecemasan" untuk akaun penting anak (seperti Apple Legacy Contact, warisan akaun WeChat).
  \item \textbf{NFC (Near Field Communication): Komunikasi jarak dekat, untuk pembayaran mudah alih, kunci pintu, perkongsian data.}: "Akta Perlindungan Maklumat Peribadi" menetapkan, proses maklumat peribadi belum cukup umur 14 tahun perlu dapat persetujuan penjaga, dan patut buat peraturan privasi khusus. Ibu bapa patut secara berkala semak polisi privasi App biasa guna anak, tutup suis "cadangan peribadi" dan "perkongsian data".
  \item \textbf{GPS (Global Positioning System): Sistem satelit untuk penentuan kedudukan.}: Akaun didaftar untuk anak patut guna e-mel/nombor telefon yang ibu bapa boleh kawal, mudah urus secara bersatu semasa belum dewasa; selepas anak dewasa boleh pindah akaun.
\end{itemize}

\begin{tcbwarn}{Wi-Fi 6/7: Standard Wi-Fi generasi baharu, kelajuan dan kestabilan lebih tinggi.}
Anak dari lahir sudah direkod - foto pertumbuhan ibu bapa pos, video taman kanak-kanak hantar, data kesihatan jam pintar. Data ini sekali muat naik sukar padam sepenuhnya. Sebelum pos fikir: sepuluh tahun lagi anak akan

\end{tcbwarn}

\begin{tcbsummary}{Windows 11: Tetapan Sistem Papar Pelbagai paparan Sambung ke paparan wayarles.}
\begin{itemize}
  \item USB-C: Port USB generasi baharu, reversible, sokong kuasa tinggi dan data pantas.
  \item PD (Power Delivery): Protokol cas pantas USB, boleh cas sehingga 240W.
  \item iOS Masa Skrin dan Android Family Link/kawalan ibu bapa jenama adalah dua garis pertahanan teknikal, guna bersama lebih stabil
  \item Mod remaja dan anti-ketagihan permainan adalah keperluan undang-undang, tetapi perlu ibu bapa aktif buka dan jaga kata laluan
  \item Literasi internet lebih penting daripada sekatan teknikal: Ajar anak kenal pasti bullying, lindungi privasi, kritikal kandungan
  \item Tabiat sihat mula dari ibu bapa - anda letak telefon, anak baru akan letak
  \item HDR (High Dynamic Range): Julat dinamik tinggi, gabung berbilang dedahan untuk gambar lebih baik.
\end{itemize}

\end{tcbsummary}
